El obrador artesano \textbf{ChocoArte} elabora y comercializa tres productos por lotes: \textbf{tabletas clásicas (T)}, \textbf{cajas de bombones (B)} y \textbf{figuras de chocolate (F)}. Los lotes pueden ser fraccionados (medio lote, un cuarto de lote, etc.). Por cada lote, el obrador obtiene un beneficio neto de \textbf{25, 30 y 45 euros}, respectivamente.

Cada semana el obrador dispone de \textbf{40 kg de cacao}: un lote de tabletas o de bombones consume \textbf{1 kg}, mientras que un lote de figuras, por su mayor tamaño, consume \textbf{2 kg}. Además, la elaboración de cualquier lote (atemperado, moldeado y envasado) ocupa \textbf{2 horas de obrador}, y la disponibilidad semanal es de \textbf{60 horas}.

Por otra parte, ChocoArte tiene un compromiso con una cadena de pastelerías por el que debe suministrar en total \textbf{al menos 12 lotes semanales}, entre los tres productos.

El modelo de programación lineal que permite determinar el plan de producción semanal óptimo es el siguiente, donde $x_1$, $x_2$ y $x_3$ representan los lotes semanales de tabletas, bombones y figuras, respectivamente:

\begin{equation}
\begin{split}
\mbox{max. } z = 25x_1 + 30x_2 + 45x_3\\
s.a.:\\
x_1 + x_2 + 2x_3 \leq 40\\
2x_1 + 2x_2 + 2x_3 \leq 60\\
x_1 + x_2 + x_3 \geq 12\\
x_1,\,x_2,\,x_3\geq 0
\end{split}
\end{equation}

Se pide:

\begin{enumerate}
    \item \textbf{Formular} el problema equivalente en formato estándar y el problema correspondiente a la primera fase del método de las dos fases.
    \item Se sabe que en la solución óptima se producen bombones y figuras y que el compromiso con la cadena de pastelerías se supera con holgura, es decir, que las variables básicas son $x_2$, $x_3$ y $h_3$. \textbf{Construir}, mediante el método de la matriz completa y sin realizar iteraciones, la tabla del símplex asociada a esa base, indicando $B$, $B^{-1}$, el valor de las variables básicas y el beneficio semanal.
    \item \textbf{Justificar}, mediante el criterio del simplex, que la base anterior es efectivamente óptima, calculando los multiplicadores del símplex y los costes reducidos.
    \item \textbf{Interpretar} económicamente el precio sombra del cacao (¿interesaría comprar cacao adicional y hasta qué precio por kilogramo?) y \textbf{calcular} el rango dentro del cual puede variar la disponibilidad semanal de cacao sin que cambie la gama de productos de la solución óptima.
    \item Por limitaciones de moldes, el obrador estudia limitar la producción de figuras a un máximo de \textbf{8 lotes semanales} ($x_3 \leq 8$). \textbf{Incorporar} esta restricción a la tabla óptima y obtener el nuevo plan de producción aplicando el método de Lemke (símplex dual), interpretando el efecto sobre el beneficio.
\end{enumerate}
